\par\indent\SourceLeaf{214}%
Ein Dreiecksnetz erstreckt sich immer nur über einen recht begrenzten Teil der Erdoberfläche. Der erwähnte Zweck lässt sich daher noch besser erreichen, wenn man die allgemeine Lösung etwas weiter verallgemeinert und von \(f\upsilon=\upsilon\) zu \(f\upsilon=\upsilon+\text{Const.}\) übergeht. Offensichtlich wäre nichts gewonnen, wenn man dieser Konstanten einen reellen Wert gäbe. Denn dann würden sich \(T\) und \(t\) nur um diese Konstante unterscheiden, und damit wären allein die Ausgangspunkte der Längen verschieden. Ganz anders verhält es sich, wenn man der Konstanten einen imaginären Wert gibt. Setzt man sie \(=-i\log k\), so ergibt sich
\[
T=t,\qquad \mathrm{tang}\tfrac{1}{2}U=k\,\mathrm{tang}\tfrac{1}{2}w\mathord{.}(\frac{1+\varepsilon\cos w}{1-\varepsilon\cos w})^{\frac{1}{2}\varepsilon}
\]
\par\indent Um den zweckmäßigsten Wert von \(k\) zu bestimmen, müssen wir zuerst das Vergrößerungsverhältnis ermitteln.

In der Schreibweise der Artikel 5 und 7 gilt hier
\[
\begin{aligned}
n&=aa\sin u^2\\
N&=AA\sin U^2\\
\varphi\upsilon&=1
\end{aligned}
\]
\par\indent Also gilt
\[
\begin{aligned}
m&=\frac{A\sin U}{a\sin u}=\frac{A\sin U}{a\sin w}\mathord{.}\BareRoot{(1-\varepsilon\varepsilon\cos w^2)}\\
&=\frac{A}{a}\mathord{.}\frac{k(1-\varepsilon\varepsilon\cos w^2)^{\frac{1}{2}+\frac{1}{2}\varepsilon}}{\cos\frac{1}{2}w^2(1-\varepsilon\cos w)^{\varepsilon}+kk\sin\frac{1}{2}w^2(1+\varepsilon\cos w)^{\varepsilon}}
\end{aligned}
\]
\noindent und dieses Verhältnis hängt somit allein von der Breite ab. Die geringstmögliche Abweichung von vollkommener Ähnlichkeit erhält man, wenn man \(k\) so bestimmt, dass \(m\) an den äußersten Breiten gleich große Werte annimmt. Dadurch liegt \(m\) bei der mittleren Breite von selbst sehr nahe an seinem größten oder kleinsten Wert. Bezeichnet man die äußersten Werte von \(w\) mit \(w^0\) und \(w'\), so erhält man auf diese Weise
\[
k=\BareRoot{\frac{\displaystyle\frac{\cos\frac{1}{2}{w^0}^{2}(1-\varepsilon\cos w^0)^{\varepsilon}}{(1-\varepsilon\varepsilon\cos {w^0}^{2})^{\frac{1}{2}+\frac{1}{2}\varepsilon}}-\frac{\cos\frac{1}{2}{w'}^{2}(1-\varepsilon\cos w')^{\varepsilon}}{(1-\varepsilon\varepsilon\cos {w'}^{2})^{\frac{1}{2}+\frac{1}{2}\varepsilon}}}{\displaystyle\frac{\sin\frac{1}{2}{w'}^{2}(1+\varepsilon\cos w')^{\varepsilon}}{(1-\varepsilon\varepsilon\cos {w'}^{2})^{\frac{1}{2}+\frac{1}{2}\varepsilon}}-\frac{\sin\frac{1}{2}{w^0}^{2}(1+\varepsilon\cos w^0)^{\varepsilon}}{(1-\varepsilon\varepsilon\cos {w^0}^{2})^{\frac{1}{2}+\frac{1}{2}\varepsilon}}}}
\]
\noindent Um zu ermitteln, bei welcher Breite \(m\) seinen größten oder kleinsten Wert annimmt, verwenden wir
\SourceLeaf{215}%
\[
\begin{aligned}
\frac{\mathrm{d}m}{m}&=\mathrm{cotang}\,U\mathord{.}\mathrm{d}u-\mathrm{cotang}\,w\mathord{.}\mathrm{d}w+\frac{\varepsilon\varepsilon\cos w\mathord{.}\sin w\mathord{.}\mathrm{d}w}{1-\varepsilon\varepsilon\cos w^2}\\
\frac{\mathrm{d}U}{\sin U}&=\frac{\mathrm{d}w}{\sin w}-\frac{\varepsilon\varepsilon\sin w\mathord{.}\mathrm{d}w}{1-\varepsilon\varepsilon\cos w^2}=\frac{(1-\varepsilon\varepsilon)\mathrm{d}w}{(1-\varepsilon\varepsilon\cos w^2)\sin w}
\end{aligned}
\]
\noindent und daraus folgt
\[
\frac{\mathrm{d}m}{m}=\frac{(1-\varepsilon\varepsilon)\mathrm{d}w}{\sin w(1-\varepsilon\varepsilon\cos w^2)}\mathord{.}(\cos U-\cos w)
\]
\noindent Daraus geht hervor, dass \(m\) dort seinen größten oder kleinsten Wert annimmt, wo \(U=w\) gilt. Bezeichnet man den Wert von \(w\) an dieser Stelle mit \(W\), so erhält man
\[
k=(\frac{1-\varepsilon\cos W}{1+\varepsilon\cos W})^{\frac{1}{2}\varepsilon}\quad\text{oder}\quad \cos W=\frac{1-k^{\frac{2}{\varepsilon}}}{\varepsilon(1+k^{\frac{2}{\varepsilon}})}
\]
\noindent Daraus lässt sich \(W\) bestimmen, wenn \(k\) nach der obigen Formel berechnet ist. In der Praxis kommt es jedoch wenig darauf an, dass die Werte von \(m\) an den äußersten Breiten genau übereinstimmen. Man kann sich damit begnügen, für \(90^0-W\) ungefähr die mittlere Breite zu wählen und daraus \(k\) abzuleiten. Den allgemeinen Zusammenhang zwischen \(U\) und \(w\) gibt dann die Formel
\[
\mathrm{tang}\tfrac{1}{2}U=\mathrm{tang}\tfrac{1}{2}w\left\{\frac{(1-\varepsilon\cos W)(1+\varepsilon\cos w)}{(1+\varepsilon\cos W)(1-\varepsilon\cos w)}\right\}^{\frac{1}{2}\varepsilon}
\]
\noindent Für die eigentliche Zahlenrechnung ist es jedoch günstiger, Reihen zu verwenden. Man kann ihnen verschiedene Formen geben, auf ihre Entwicklung gehen wir hier aber nicht ein.

Man sieht leicht, dass für \(w<W\) die Ungleichung \(U>w\) gilt, sodass \(\cos U-\cos w\) und damit auch \(\frac{\mathrm{d}m}{\mathrm{d}w}\) negativ sind. Für \(w>W\) gilt \(U<w\), sodass \(\frac{\mathrm{d}m}{\mathrm{d}w}\) positiv ist. Daher ist klar, dass für \(w=U=W\) der Wert von \(m\) stets ein Minimum annimmt, und zwar
\[
=\frac{A}{a}\BareRoot{(1-\varepsilon\varepsilon\cos W^2)}
\]
\noindent Wählt man also den Kugelradius \(A=\frac{a}{\BareRoot{(1-\varepsilon\varepsilon\cos W^2)}}\), dann ist die Abbildung unendlich kleiner Teile des Ellipsoids bei der Breite \(90^0-W\) dem Urbild ähnlich und sogar gleich. Bei anderen Breiten ist sie größer.

Man kann den Logarithmus von \(m\) mit Vorteil in eine Reihe nach aufeinanderfolgenden Potenzen von \(\cos U-\cos W\) entwickeln. Ihre ersten, für die Praxis ausreichenden Glieder sind
\SourceLeaf{216}%
\[
\begin{aligned}
\log\mathrm{hyp}.\,m&=\log\{\frac{A}{a}\BareRoot{(1-\varepsilon\varepsilon\cos W^2)}\}+\frac{\varepsilon\varepsilon}{2(1-\varepsilon\varepsilon)}\mathord{.}(\cos U-\cos W)^2\\
&\qquad-\frac{2\varepsilon^4\cos W}{3(1-\varepsilon\varepsilon)^2}\mathord{.}(\cos U-\cos W)^3\ .\ .\ .\ .
\end{aligned}
\]
\par\indent Überträgt man also zum Beispiel die Dänische Monarchie zwischen den Breitengrenzen \(53^0\) und \(58^0\) auf diese Weise auf die Kugeloberfläche und setzt \(W=34^0\,30'\), so wird die Abbildung bei der Abplattung \(\frac{1}{303}\) an den Grenzen in ihren Längenmaßen nur um \(\frac{1}{530000}\) vergrößert.

Wir müssen uns hier damit begnügen, \textit{eine} Anwendung der Übertragung von Figuren in der höheren Geodäsie kurz anzudeuten, und eine angemessene Ausarbeitung für einen anderen Ort aufsparen.

